\documentclass[11pt]{article}
\usepackage[margin=1in]{geometry}
\usepackage{amsmath,amssymb,amsthm}
\usepackage[hidelinks]{hyperref}

\newtheorem*{theorem}{Literature Answer}

\title{Literature Answer to Problem 1.13 of arXiv:math/0508650}
\author{FA/Banach run packet}
\date{2026-06-18}

\begin{document}
\maketitle

\section*{Status}

This is a literature-already-answered packet. It records that Problem 1.13
from Dilworth--Odell--Sari \cite{DOS} was answered in the later paper of
Dodos \cite{Dodos}. It is not an original proof from this run.

\section*{Original Problem}

Problem 1.13 appears in Section 1 of \cite{DOS}. In the notation of that
paper, it asks:

\begin{quote}
If \(SP_w(X)\) is uncountable, must there exist
\(\{(\tilde x_i^\alpha)_{i=1}^\infty:\alpha<\omega_1\}\subseteq SP_w(X)\)
which is either strictly increasing with respect to \(\alpha\), strictly
decreasing, or consists of mutually incomparable elements?
\end{quote}

\section*{Supporting Answer}

Dodos \cite{Dodos} explicitly introduces the same question as the problem
posed by the authors of \cite{DOS}, citing it as \cite[Problem 1.13]{DOS}.
The abstract states that, for every separable Banach space \(X\), either
\(SP_w(X)\) is countable or \(SP_w(X)\) contains an antichain of cardinality
continuum, and says that this answers a question of Dilworth, Odell, and Sari.

\begin{theorem}
Problem 1.13 of \cite{DOS} has a positive answer for separable Banach spaces.
If \(SP_w(X)\) is uncountable, then \(SP_w(X)\) contains a family of
cardinality continuum whose elements are pairwise incomparable in the
domination order.
\end{theorem}

\begin{proof}[Literature identification]
Theorem 1 of \cite{Dodos} states that, whenever \(SP_w(X)\) is uncountable,
there is a Cantor-tree family \((x_t)_{t\in 2^{<\mathbb N}}\) in \(X\) such
that different branches generate spreading models incomparable under
domination. Corollary 1 of \cite{Dodos} records the direct consequence:
\(SP_w(X)\) contains an antichain of size continuum.

Since continuum has cardinality at least \(\omega_1\), this continuum-sized
antichain contains an \(\omega_1\)-indexed subfamily. That subfamily satisfies
the third alternative in Problem 1.13 of \cite{DOS}, namely mutual
incomparability. Thus the later paper answers the original problem.
\end{proof}

\section*{Scope Limitations}

This packet answers only Problem 1.13 of \cite{DOS}. It does not address the
countable-spreading-model Problem 1.6, nor the broader set-theoretic
semilattice Problem 1.14 in the same source paper.

\section*{Verification Notes}

The source and supporting papers are distinct. The supporting paper explicitly
cites \cite[Problem 1.13]{DOS}, so this belongs in
\texttt{literature\_already\_answered} rather than
\texttt{literature\_implied\_answers}. The local duplicate search covered
arXiv id \texttt{0508650}, title keywords \texttt{Lattice structures and
spreading models}, and phrases \texttt{Problem 1.13}, \texttt{uncountable
SP\_w}, and \texttt{antichain of spreading models}.

\begin{thebibliography}{9}

\bibitem{DOS}
S. J. Dilworth, E. Odell, and B. Sari,
\emph{Lattice structures and spreading models},
arXiv:math/0508650.

\bibitem{Dodos}
P. Dodos,
\emph{On antichains of spreading models of Banach spaces},
arXiv:0805.2038.

\end{thebibliography}

\end{document}
